%% ======================================================================
%% 02_context.tex -- 1. Hardware in context
%% ======================================================================
\section{Hardware in context}
\label{sec:context}

Balancing systems, such as the inverted pendulum, ball and beam, ball and plate, and multirotor attitude testbeds, are standard benchmarks in control education and research because they are open-loop unstable, inexpensive, and visually intuitive \cite{astrom2008,hauser1992,awtar2002}. Open versions published in HardwareX include a multi-arm pendulum on a cart \cite{kaheman2023} and a low-cost video-tracked double pendulum \cite{myers2020}, alongside open educational mobile robots such as ROBOTONT \cite{raudmae2023} and PlatROB \cite{balbuena2026}, and educational manipulators and heads such as EduSCARA \cite{clark2026} and BEATRIX \cite{al2024}. These platforms study one balancing problem at a time on a fixed base.

In field robotics, the relevant scenario is the opposite: a payload that must stay level rides on a vehicle that accelerates, turns, and climbs slopes, and the vehicle itself may carry an actively stabilized stage. Examples include camera gimbals on ground vehicles, stabilized trays for fragile objects \cite{selvaggio2023}, and aerial vehicles balancing an unstable payload \cite{hehn2011}. In these systems, disturbances generated by the base propagate through each stage; every stabilizer attenuates part of the disturbance while adding its own dynamics, delays, and saturations. Prior projects in our university course addressed a single stabilized stage (a two-axis gimbal on a differential robot tracking the ball with a resistive touchscreen or a camera). To our knowledge, no documented open platform allows disturbance propagation through \emph{two} actively controlled stages---one of them an aerial vehicle---to be measured with synchronized logs \emph{and} predicted with a validated model.

\paragraph{Why a drone instead of a gimbal}
A natural question is why the intermediate stage is an octocopter rather than a motorized two-axis gimbal, which would also level the plate. Section~\ref{sec:porquedron} gives the quantitative answer. A gimbal can only \emph{tilt} the payload; the octocopter levels itself with the same control law while also \emph{supporting its own weight}. That vertical force is the only force on the platform that bypasses the wheel--ground contact, changing the base's limits on a ramp: with half of the module's weight carried by the rotors, the static uphill tip-over angle rises from \SI{21.5}{\degree} (measured) to a predicted \SI{27.7}{\degree}, the downhill tip-over angle rises from \SI{30.5}{\degree} (measured) to a predicted \SI{38.2}{\degree}, and the acceleration margin on a \SI{20}{\degree} slope increases by a factor of 3.5. At the same time, the drone is the most demanding intermediate stage (an inverted pendulum with actuator lag and limited torque authority), making the platform a safe, measurable stepping stone toward free-flight aerial manipulation, our group's research focus.

AeroTray fills this gap using off-the-shelf components. Its contributions are:
\begin{enumerate}[label=(\roman*)]
  \item a reproducible three-tier architecture (ground base, octocopter on a passive spherical joint with a vertical guide, and ball-and-plate stage) with documented interfaces across three independent controllers (myRIO, ESP32, ESP32-S3) and a Raspberry Pi supervisor, supported by isolated power domains and an emergency stop (Section~\ref{sec:description});
  \item a symbolically verified first-principles kinematic and dynamic model of the full chain, from wheel velocities to ball acceleration, covering the nonholonomic base on a slope, the octocopter as an inverted pendulum on an accelerating pivot, eight-rotor control allocation, and exact equations of a ball rolling on a translating and rotating plate (Section~\ref{sec:model} and \ref{app:derivations});
  \item a set of model-based results explaining prototype behavior: center-of-mass identification of the aerial module from hover telemetry, quantitative justification of the drone on a ramp (tip-over, torque, and traction versus supported weight fraction), thrust margin and endurance versus battery voltage, an attitude-loop stability map, a torque-authority region of attraction, the delay-induced stability limit of the ball loop, transmissibility from base accelerations to the ball, and static and dynamic base tip-over limits (Sections~\ref{sec:porquedron}--\ref{sec:modelval});
  \item open firmware for an octocopter attitude controller with guided-flight mode, safety interlocks, and onboard logging; a \SI{187}{fps} vision-based ball controller on a Raspberry Pi 3B+; and quantitative validation including failed iterations, comparison with the Isaac Lab model used during design \cite{mittal2023,schulman2017}, telemetry-synchronized flight videos, and all raw data (Section~\ref{sec:validation}).
\end{enumerate}
The platform is designed to teach integration, modeling, and validation in mechatronic design, and to serve as a testbed for cascaded, learning-based, and disturbance-feedforward controllers under base-induced perturbations.
